\section{Mathematical foundations}
\label{sec:math}
This section collects the statements that justify the engine, the statements
that show where it goes wrong, and the two identities used by the corrected
version. All three reviews prove variants of
Propositions~\ref{prop:gcd}--\ref{prop:rat}; the power-identity conditions of
Proposition~\ref{prop:power} and the quadratic-surd formula of
Proposition~\ref{prop:surd} are stated in \RIII\ in a weaker form.

\subsection{Why the GCD step is meaningful}
\begin{proposition}\label{prop:gcd}
Let $\rho\ne0$ and $m\ne0$ be algebraic numbers, $q\ge2$, and let $K\supseteq\Q$
be the field generated by the coefficients used by \code{PolynomialGCD} with
\code{Extension -> Automatic} (it contains $m\rho$). Let $P_m$ be the minimal
polynomial over $\Q$ of $\alpha=(m\rho)^{1/q}$ and $G_m=\gcd_{K[X]}(P_m,X^q-m\rho)$.
Then $1\le\deg G_m\le q$, and every root of $G_m$ is a $q$-th root of $m\rho$.
\end{proposition}
\begin{proof}
$\alpha$ is a common root, so the minimal polynomial of $\alpha$ over $K$ divides
both polynomials and hence their GCD; thus $\deg G_m\ge1$. $G_m$ divides the
degree-$q$ binomial, which gives the upper bound and the root property.
\end{proof}
A GCD of degree $e<q$ is an equation of degree $e$ over $K$ for some $q$-th
root of $m\rho$. In the best case $e=1$ and that root is an element of $K$
written without the outer radical. Nothing in the proposition says that
solving $G_m$ yields a \emph{simpler radical expression}: the roots may involve
\code{Root} objects or nested radicals of their own (experiments Q01--Q08 show
candidates containing $\sqrt{9-6\sqrt2}$, which is itself denestable).

\subsection{Why the roots-of-unity orbit is right, and why the filter is wrong}
\label{sec:branch}
\begin{proposition}\label{prop:orbit}
Let $\beta\ne0$, $q\ge2$, $\rho=\beta^q$, $m\ne0$, and let $r$ be any root of
$G_m$. Then $\{\,r\,m^{-1/q}\zeta_q^{\,k}: 0\le k<q\,\}$, $\zeta_q=e^{2\pi i/q}$, is
exactly the set of $q$-th roots of $\rho$; in particular it contains $\beta$.
\end{proposition}
\begin{proof}
$r^q=m\rho$ and $(m^{1/q})^q=m$ for any choice of $m^{1/q}$, so
$(r/m^{1/q})^q=\rho$. The $q$ numbers are distinct since $\rho\ne0$, and
$X^q-\rho$ has exactly $q$ roots.
\end{proof}
No identity of the form $(m\rho)^{1/q}=m^{1/q}\rho^{1/q}$ is used. The code
builds precisely this orbit (\src{121--124}); the orbit is a strength of the
program. The defect is the selection:
\begin{lstlisting}
Select[possibilities, PossibleZeroQ[# - radicand^outerpower] &]
\end{lstlisting}
compares against $\rho^{1/q}=(\beta^q)^{1/q}$, not against $\beta$.
\begin{proposition}\label{prop:principal}
For $\beta=re^{i\theta}\ne0$ with $\theta=\Arg\beta\in(-\pi,\pi]$ and an
integer $q\ge2$,
\[
(\beta^q)^{1/q}=\beta\,e^{-2\pi ik/q},\qquad k\in\Z\text{ with }q\theta-2\pi k\in(-\pi,\pi].
\]
Hence $(\beta^q)^{1/q}=\beta$ if and only if $-\pi/q<\Arg\beta\le\pi/q$.
\end{proposition}
\begin{proof}
$\Arg(\beta^q)=q\theta-2\pi k$ by definition of the principal argument; take the
principal $q$-th root.
\end{proof}
Every input to the engine that is itself a principal root
($\beta=\rho^{1/q}$ literally) lies in the sector, which is why the classical
examples work. The wrapper, however, sends products of markers and powers
$\text{rad}^{p/q}$ with $p>1$ to the engine; those need not lie in the sector.
Example: $\beta=(-1+2i\sqrt2)^{3/2}=(1+i\sqrt2)^3=-5+i\sqrt2$ has
$\Arg\beta>\pi/2$, so $(\beta^2)^{1/2}=5-i\sqrt2=-\beta$, and the program returns
$5-i\sqrt2$ (experiment C02). The randomized family F5 shows that this is the
typical, not the exceptional, outcome for such inputs.

\subsection{The rationalizer is sound}
\begin{proposition}\label{prop:rat}
Let $d\ne0$ be algebraic with minimal polynomial $P\in\Q[X]$ (any nonzero
rational multiple, not necessarily monic), and let $P(X)=(X-d)Q(X)$. Then
$P(0)\ne0$ and $1/d=-Q(0)/P(0)$.
\end{proposition}
\begin{proof}
$P$ is irreducible of degree $\ge1$ with root $d\ne0$, so $P(0)\ne0$. Evaluating
$P(X)=(X-d)Q(X)$ at $0$ gives $P(0)=-dQ(0)$.
\end{proof}
This is exactly \src{447--450}. Experiments R01--R09 confirm the helper on
quadratic, quartic, complex, cubic, and nested denominators.

\subsection{When power identities are valid}
\begin{proposition}\label{prop:power}
For nonzero complex $a,F,S$ and real exponents:
\begin{enumerate}
\item $(a^b)^c=a^{bc}$ holds whenever $c\in\Z$, or $b\in\R$ with $-1<b\le1$.
\item Let $p=n/q$ in lowest terms with $q\ge2$. Then $(FS)^p=F^pS^p$ if and only
      if $\Arg F+\Arg S\in(-\pi,\pi]$. In particular it holds whenever $F>0$ or
      $S>0$.
\end{enumerate}
\end{proposition}
\begin{proof}
(1) If $c\in\Z$, $(a^b)^c=\exp(c\log(a^b))$ and $\exp(c\cdot 2\pi i k)=1$
absorbs the branch. If $-1<b\le1$ then $\operatorname{Im}(b\log a)=b\Arg a\in(-\pi,\pi]$,
so $\log(a^b)=b\log a$ exactly and $(a^b)^c=\exp(cb\log a)=a^{bc}$.
(2) $\log(FS)=\log F+\log S-2\pi ik$ with $k\in\{-1,0,1\}$ chosen so that the
imaginary part lies in $(-\pi,\pi]$; $k=0$ iff $\Arg F+\Arg S\in(-\pi,\pi]$. Then
$(FS)^p=F^pS^p\exp(-2\pi ikp)$, and $\exp(-2\pi ikn/q)=1$ iff $q\mid kn$, i.e.\
iff $k=0$ since $\gcd(n,q)=1$ and $|k|\le1<q$.
\end{proof}
The custom factorizer applies (1) with rational $c$ and unrestricted $b$
(\src{527--531}) and (2) with rational $p$ and unrestricted $F$
(\src{571--583}). Counterexamples for (1): $\sqrt{z^2}\mapsto z$ (E01).
For (2): $\sqrt{-x-y}\mapsto i\sqrt{x+y}$ (E03), and, entirely numerically,
$F=-1$, $S=\sqrt2+(-1)^{1/3}$ with $\Arg F+\Arg S=\pi+\pi/9>\pi$, where the
program returns $i\sqrt S=-\sqrt{-S}$ (J01). The corrected factorizer applies
(1) only under the stated condition and (2) only when $p\in\Z$, $F>0$ or $S>0$.

\subsection{Quadratic surds}
\begin{proposition}\label{prop:surd}
Let $a,b,c\in\Q$ with $a>0$, $b\ne0$, $c>0$ not a square, and suppose
$\Delta=a^2-b^2c=\delta^2$ with $\delta\in\Q_{\ge0}$. Put $u=(a+\delta)/2$ and
$v=(a-\delta)/2$. Then $u\ge v\ge0$ and
\[
\sqrt{a+b\sqrt c}=\sqrt u+\operatorname{sgn}(b)\sqrt v ,
\]
where both sides are principal square roots. If $a<0$ and $a+b\sqrt c<0$ then
$\sqrt{a+b\sqrt c}=i\sqrt{-a-b\sqrt c}$ and the formula applies to $-a,-b$.
\end{proposition}
\begin{proof}
$a\ge\delta$ because $a^2=\Delta+b^2c\ge\Delta$, so $v\ge0$. The square of the right
side is $u+v+2\operatorname{sgn}(b)\sqrt{uv}=a+\operatorname{sgn}(b)\sqrt{a^2-\Delta}
=a+\operatorname{sgn}(b)|b|\sqrt c=a+b\sqrt c$, and the right side is
$\ge\sqrt u-\sqrt v\ge0$, so it is the principal root. The negative case is
$\sqrt{-w}=i\sqrt w$ for $w>0$.
\end{proof}
This is the classical criterion (one half of Landau's theorem for quadratic
surds); the corrected version uses it both as a fast path and to clean inner
surds that appear inside candidates, e.g.\ $\sqrt{9-6\sqrt2}=\sqrt6-\sqrt3$.

\subsection{How many multipliers a batch contains}
\begin{proposition}\label{prop:count}
If $t$ distinct primes are selected, the batch built from
$\operatorname{Divisors}\big((\prod_{i\le t}p_i)^{q-1}\big)$ contains exactly $q^t$
elements before the unit is dropped.
\end{proposition}
\begin{proof}
Each divisor is $\prod p_i^{e_i}$ with independent $0\le e_i\le q-1$.
\end{proof}
The code selects $t=\min(\#\text{primes},\max(5,\lceil\log1000/\log q\rceil))$
(\src{190--194}), so the nominal cap of 1000 is not a cap: $t=10$ for $q=2$
gives 1023 candidates, and the forced minimum $t=5$ gives $q^5-1$ for large
$q$ ($99\,999$ for $q=10$; experiment N06 materializes exactly that list). The
corrected version chooses the largest $t$ with $q^t-1\le\text{cap}$ in integer
arithmetic.
